\documentclass[12pt]{article}
\usepackage{amsmath,amsfonts,fullpage,amsthm,graphicx,verbatim}

\newtheorem*{theorem}{Theorem}
\newtheorem*{fact}{Fact}
\newtheorem*{goal}{Goal}
\newtheorem*{idea}{Idea}
\newtheorem*{defn}{Definition}
\newtheorem*{ex}{Example}

%Style section
%\setlength{\textheight}{10in}
%\setlength{\textwidth}{7in}
%\hoffset=-0.75in
 \pagestyle{plain}
% \setlength{\topmargin}{0in}
% \setlength{\headsep}{0in}
% \setlength{\oddsidemargin}{0in}
% \setlength{\evensidemargin}{0in}

\begin{document}


\pagestyle{empty} %use this to get rid of page numbers

\noindent
{Math 166 \hskip 1 truein  Names:\ \hrulefill}

\noindent
%\vskip 0.3truein
%\bigskip
%\noindent
{Work on work}

\begin{goal} \rm
Use the integral, our infinite adding machine, to find the work required to do the following.
\end{goal}
%\smallskip


\begin{enumerate}
%\setcounter{enumi}{-1}

\item An underground hemispherical dome with base a circle of radius 20 m is filled with mud of density 1906 kg/$m^3$. (Note the dome  is oriented so the flat part is farthest below ground.) Find the work done against gravity to pump enough mud out of the dome so that the level of the mud in the dome is 2 m after pumping. The dome is buried 1 meter below ground and the mud travels up the pipe and is dispensed on top of the ground. Draw a relevant picture in the coordinate axes and explain each step of your work.

%\begin{comment}
\begin{enumerate}
\item First draw a picture of the dome in the coordinate axes. 
%with the $y$-axis as the vertical axis and the origin at the bottom of the tank. 
Draw a horizontal layer at height $y$.
\vfill\vfill
\item Next, find the area $A(y)$ of the layer at height $y$ as a function of $y$.
\vfill
\item Write down an expression for the volume of this layer if it has thickness $\Delta y$.
\vfill
\item What is the mass of this layer?
\vfill
\pagebreak
\item What is the force due to gravity on this layer? 
\vfill
\item What is the vertical distance that this layer is lifted?
\vfill
\item Use your answers to (e) and (f) to find the work needed to pump the mud in this layer out of the spout.
\vfill
\item Integrate this expression to solve the problem.
\vfill
\end{enumerate}
%\end{comment}

\pagebreak

\item A 35 meter rope whose mass is 105 kg is hanging from the top of a building 35
meters high. A 250 kg piano is sitting on the ground attached to the lower end of the rope. How much work
is needed to reel in the rope and the piano so that the piano is at a level that is 10 meters below the top of the building? Draw a relevant picture in the coordinate axes and explain each step of your work.
\pagebreak
\item A tank of water is 15 feet long and has a cross section in the shape of an equilateral triangle with sides 2 feet long (point of the triangle points directly down). The tank is filled with water to a depth of 9 inches. Determine the amount of work needed to pump all of the water to the top of the tank. Assume that the weight of the water is $62\; lb/ft^3$. 
\begin{comment}
\begin{enumerate}
\item Draw a picture of the crane and wrecking ball in the coordinate axes, with the $y$-axis as the vertical axis and the origin at the initial position of the wrecking ball on the ground. 
\bigskip
\bigskip
\bigskip
\bigskip
\bigskip
\bigskip
\bigskip
\bigskip
\item First, find the work needed to lift the wrecking ball 12 m, neglecting the work needed to lift the cable. (Hint: no integral is needed for this part. Use Work = Force $\times$ Distance and Force = Mass $\times$ g.)
\bigskip
\bigskip
\bigskip
\item Now consider a segment of cable at height $y$ of length $\Delta y$.
What is the mass of this segment?
\bigskip
\bigskip
\bigskip
\item What is the force due to gravity on this segment? 
\bigskip
\bigskip
\bigskip
\item What is the vertical distance that this segment is lifted?
\bigskip
\bigskip
\bigskip
\item Multiply (d) and (e) to find the work needed to lift this segment.
\bigskip
\bigskip
\bigskip
\item Integrate this expression to find the work to lift the cable (evaluate the integral only if you have time).
\bigskip
\bigskip
\bigskip
\item Add together (a) and (g) to find the work to lift the cable and wrecking ball.
\end{enumerate}
\end{comment}

\end{enumerate} 

\end{document}
